%%%%%%%%%%%%%%%%%%%%%%%%%%%%%%%%%%%%%%%%%%%%%%%%%%%%%%%%%%%%%%%%%%%%%%
% Lớp 10 | Phương pháp tọa độ trong mặt phẳng (𝑂𝑥𝑦) | Ba đường conic trong mặt phẳng tọa độ
% Dạng 0H9?5-6: Bài toán điểm trên hypebol
% Số câu: 2  (sinh tự động bởi tach_cau.py; từ nay thầy sửa trực tiếp file này)
%%%%%%%%%%%%%%%%%%%%%%%%%%%%%%%%%%%%%%%%%%%%%%%%%%%%%%%%%%%%%%%%%%%%%%

%%%%%%%%%%%%%%%% ĐÚNG/SAI %%%%%%%%%%%%%%%%

% nguồn: MaHoaNH_DungSai_Lop10
\begin{ex}%[0H9H5-6]%[0H9V5-4]%[ma:94fb2dc0]
	 Cho hypebol $(H): \dfrac{x^2}{16}-\dfrac{y^2}{9}=1$.
	 \choiceTF
		{Tiêu cự  của hypebol bằng $5$}
		{Điểm $A(4; 3) \in(H)$} 
		{\True Tiêu điểm của hypebol là  $F_1(-5; 0), F_2(5; 0)$}
		{\True Trên $(H)$ có 4 điểm $M$ sao cho $MF_1 \perp MF_2$} 
	\loigiai{\begin{itemchoice}
		\itemch Ta có $a^2=16\Rightarrow a=4, b^2=9\Rightarrow b=3;  c^2=a^2+b^2=16+9=25\Rightarrow c=5$.
		\itemch Thay tọa độ điểm $A(4;  3)$ vào phương trình $\dfrac{x^2}{16}-\dfrac{y^2}{9}=1$, ta được $0=1$ (vô lý)
		\itemch Tiêu điểm: $F_1(-5;  0), F_2(5; 0)$
		\itemch Gọi $M\left(x_M; y_M\right) \in(H)$ và $MF_1 \perp  MF_2$ nên \allowdisplaybreaks
		\begin{eqnarray*}
			 &&MF_1^2+MF_2^2=F_1 F_2^2\\
			&\Leftrightarrow&\left(a+\dfrac{c}{a} x_M\right)^2+\left(a-\dfrac{c}{a} x_M\right)^2=(2c)^2\\ 
			&\Leftrightarrow&\left(4+\dfrac{5}{4}  x_M\right)^2+\left(4-\dfrac{5}{4} x_M\right)^2=100\\
			&\Leftrightarrow& x_M^2=\dfrac{544}{25}\\ &\Leftrightarrow&\hoac{&x_M=-\dfrac{4\sqrt{34}}{5} \Rightarrow y_M=\pm  \dfrac{9}{5} \Rightarrow M_1\left(-\dfrac{4\sqrt{34}}{5};-\dfrac{9}{5}\right), M_2\left(-\dfrac{4\sqrt{34}}{5}; \dfrac{9}{5}\right) \\
			&x_M=\dfrac{4\sqrt{34}}{5} \Rightarrow y_M=\pm  \dfrac{9}{5} \Rightarrow M_3\left(\dfrac{4\sqrt{34}}{5};-\dfrac{9}{5}\right), M_4\left(\dfrac{4\sqrt{34}}{5}; \dfrac{9}{5}\right)
			 }
		\end{eqnarray*}
		Có $4$ điểm thỏa  mãn yêu cầu:
		$M_1\left(-\dfrac{4\sqrt{34}}{5}; \dfrac{9}{5}\right)$, $M_2\left(-\dfrac{4\sqrt{34}}{5};-\dfrac{9}{5}\right)$, $M_3\left(\dfrac{4\sqrt{34}}{5};-\dfrac{9}{5}\right)$, $M_4\left(\dfrac{4\sqrt{34}}{5}; \dfrac{9}{5}\right)$. 
\end{itemchoice}}
\end{ex}


%%%%%%%%%%%%%%%% TRẢ LỜI NGẮN %%%%%%%%%%%%%%%%

% nguồn: MaHoaNH_TraLoiNgan_Lop10
\begin{bt}%[0H9H5-6]%[0T9B3-2]%[ma:0d93d3ca]
	Trong mặt phẳng tọa độ $ Oxy $, cho hai điểm $ A(-1;5) $, $ B(3;2) $. Biết phương trình $ \dfrac{x^2}{a^2}-\dfrac{y^2}{b^2}=1$ là phương  trình chính tắc của một hypebol $(H)$ có độ dài trục ảo bằng $ 2AB $ và có tiêu điểm $ F_1(-10;0) $. Tính $a^2+2b^2$.
	 \shortans{125}
	\loigiai{ 
		Độ dài trục ảo bằng $2b=2AB \Rightarrow b=AB=5$.\\ 
		Tiêu điểm hypebol $ F_1(-10;0)  \Rightarrow c=10 $.\\
		Phương trình  chính tắc của hypebol $ (H) \colon \dfrac{x^2}{75}-\dfrac{y^2}{25}=1$.\\
		Vậy  $a^2+2b^2=125$.
	} 
\end{bt}
